\section{Local families and fixed representatives}
\label{sec:local-families}

We first record the local information needed for the numerical comparison
and for the eventual ambient trivialization.  Throughout this section,
$F(x,t)=f_t(x)$ is a holomorphic family in three spatial variables over a
disk $\Delta$, with $f_t(0)=0$, an isolated critical point at $0$, and
$\mu(f_t,0)=\mu>0$ for every $t\in\Delta$.  Write
$X=F^{-1}(0)$, $X_t=f_t^{-1}(0)$, and
$\Gamma=\{0\}\times\Delta$.  Representatives and the parameter disk may
always be shrunk about the central point.

\begin{lemma}[Local family properties]
\label{lem:local-family-properties}
There are representatives on which the spatial critical locus of $F$ is
exactly $\Gamma$.  On these representatives $X\to\Delta$ is flat, its
fibers are normal Gorenstein surfaces, and $X$ is a normal Gorenstein
threefold, smooth over $\Delta$ away from $\Gamma$.  The same statements
hold after a nonconstant holomorphic base change between smooth disks.
\end{lemma}

\begin{proof}
Consider the scheme defined by the three spatial partial derivatives.
Its central fiber is supported at $0$ and has length $\mu$.  The local
finiteness theorem for analytic maps gives a finite representative over
the disk.  In the regular local ring $\mathbb C\{x_1,x_2,x_3,t\}$, the
three partials together with $t$ form a system of parameters.  They
therefore form a regular sequence.  The quotient by the three partials
is Cohen--Macaulay of dimension one, and $t$ is a nonzerodivisor in it.
It is consequently finite flat over the disk, of degree $\mu$.
In every nearby fiber its local length at $0$ is already $\mu$; there
can be no further points.  This is conservation of the spatial-gradient
number; compare \cite[Proposition~3.1]{Denkowski2018}.

Each surface fiber is a hypersurface, hence Gorenstein and
Cohen--Macaulay.  Its isolated singularity has codimension two, so
Serre's criterion gives normality.  Similarly, the total hypersurface
is Gorenstein and Cohen--Macaulay; it is smooth off $\Gamma$, which has
codimension two in $X$, and is therefore normal.  Since no fiber of the
ambient projection is contained in $X$, its parameter is a
nonzerodivisor on $X$.  Torsion-freeness over the regular one-dimensional
base gives flatness.  After a disk base change the equation is again a
relative hypersurface, with the same individual surface germs and the
same spatial critical locus.  These arguments apply again.
\end{proof}

Use $\chi_c$ for Euler characteristic with compact supports; on a compact
space it agrees with ordinary Euler characteristic $\chi$.

\begin{lemma}[Euler characteristic in a fixed ball]
\label{lem:fixed-ball-euler}
There are a sufficiently small closed Euclidean ball $B$ about $0$, a
smaller parameter disk, and $\eta>0$ such that the following statements
hold for every parameter in that disk:
\begin{enumerate}
\item $0$ is the only spatial critical point in $B$;
\item $f_t^{-1}(c)$ meets $\partial B$ transversely whenever $|c|<\eta$;
\item
\[
 \chi(X_t\cap B)=1,
 \qquad
 \chi_c\bigl((X_t\cap B)\setminus\{0\}\bigr)=0.
\]
\end{enumerate}
The Euler equalities also hold if the radius of $B$ varies slightly
within a sufficiently small interval about the chosen radius.
\end{lemma}

\begin{proof}
Choose $B$ inside the representative of
Lemma~\ref{lem:local-family-properties}, with radius a Milnor radius for
$f_0$.  Transversality at the compact set
$f_0^{-1}(0)\cap\partial B$ persists for small $t$ and small $c$.
Points elsewhere on the sphere do not enter these small levels, after
shrinking, by compactness.  This proves the first two assertions.

Fix $c$ with $0<|c|<\eta$.  The spaces $f_t^{-1}(c)\cap B$ form a proper
smooth family with boundary over the smaller parameter disk.  Smooth
lifts tangent to the boundary give the proper-submersion theorem in
this setting.  Their Euler characteristic is therefore that of the
central Milnor fiber, namely $1+\mu$; see
\cite[Theorems~6.5 and~7.2]{Milnor1968} and, for the analytic
Jacobian-algebra formulation, \cite[Section~1, Theorem~1.5]{LeRamanujam1976}.

For a fixed $t$, choose an inner Milnor ball for $f_t$, with closure in
the interior of $B$.  On the closed region between the two spheres,
$f_t$ has no critical point, and both boundary restrictions are
submersive over a sufficiently small disk of values.  Thus the shell
has the same Euler contribution at value $0$ and at sufficiently small
nonzero values.  In the inner ball, passing from value $0$ to a small
nonzero value replaces a contractible singular fiber by a Milnor fiber
of Euler characteristic $1+\mu$.  Additivity of $\chi_c$, or a
decomposition using the inner boundary and the two complementary
pieces, therefore gives
\[
 \chi(f_t^{-1}(c)\cap B)-\chi(X_t\cap B)=\mu.
\]
The inner ball and the admissible nonzero value may depend on $t$;
the preceding computation of the outer smooth level applies to every
such value.  It follows that $\chi(X_t\cap B)=1$.  Removing its one
origin point gives the second equality.  Finally, transversality
persists for radii in a small interval around the chosen radius, and
the identical argument applies there.
\end{proof}

Our algebraic reduction will concern the numerical invariant of each
individual surface germ.  The following elementary form of finite
determinacy supplies polynomial representatives for that purpose.

\begin{lemma}[Individual finite determinacy]
\label{lem:individual-jets}
Let $f\in\mathbb C\{x_1,x_2,x_3\}$ have an isolated critical point at $0$
and Milnor number $\mu$, and write $\mathfrak m=(x_1,x_2,x_3)$.
For every $b\in\mathfrak m^{\mu+2}$ there is a biholomorphic germ $h$
fixing $0$ such that $(f+b)\circ h=f$.  In particular, if
$N\geq\mu+1$, then $f$ is analytically right-equivalent to its Taylor
polynomial $j^Nf$.
\end{lemma}

\begin{proof}
Let $\mathfrak m=(x_1,x_2,x_3)$ and $J(f)=(f_{x_1},f_{x_2},f_{x_3})$.
The local algebra $\mathbb C\{x\}/J(f)$ has length $\mu$, so its maximal
ideal has $\mu$th power zero, by the descending filtration and
Nakayama's lemma.  Hence $\mathfrak m^\mu\subset J(f)$.
For the given $b\in\mathfrak m^{\mu+2}$ we have
\[
 b\in\mathfrak m^{\mu+2}\subset\mathfrak m^2J(f),
 \qquad
 b_{x_i}\in\mathfrak m^{\mu+1}\subset\mathfrak m J(f).
\]
Write $b=a^{\mathsf T}g$, where $g=(f_{x_i})_i$ is a column vector and
the entries of $a$ lie in $\mathfrak m^2$.  There is a matrix $C$ with
entries in $\mathfrak m$ such that the gradient column of $f+sb$ is
$(I+sC)g$.  On a sufficiently small neighborhood this matrix is
invertible simultaneously for $0\leq s\leq1$.  The holomorphic spatial
vector field
\[
 w_s=-(I+sC)^{-\mathsf T}a
\]
vanishes to second order and satisfies $d(f+sb)(w_s)=-b$.
Its time-dependent flow $h_s$, initially the identity, exists as a
holomorphic germ throughout $0\leq s\leq1$, after shrinking the spatial
neighborhood.  This follows from a uniform quadratic bound on $w_s$;
the reverse equation supplies the inverse germ.  The chain rule gives
$\frac{d}{ds}((f+sb)\circ h_s)=0$, and therefore $h=h_1$ has the
required property.  Taking $b=j^Nf-f$ proves the assertion about Taylor
polynomials.
\end{proof}

\subsection{Birational conventions}
A pair $(Z,B)$ consists of a normal complex variety and an effective rational
Weil divisor $B$ such that $K_Z+B$ is rational Cartier.  For a prime divisor
$E$ on a birational model, its log discrepancy is one plus the coefficient
of $E$ in the difference between the canonical divisor upstairs and the
pullback of $K_Z+B$, using compatible canonical divisors.  The pair is
\emph{log canonical}, abbreviated lc, if all log discrepancies are nonnegative.
For $B=0$ we say that $Z$ is log canonical.  A pair with boundary coefficients
at most one is \emph{divisorial log terminal}, abbreviated dlt, if it has a
log resolution on which every exceptional divisor has positive log
discrepancy.  An lc center is the image of a divisor with log discrepancy zero.
For an lc pair, the equivalent dlt criterion is that a simple
normal crossing open set contains the generic points of all lc centers.
We use the usual normality and adjunction properties of its coefficient-one strata
\cite[11.5--11.15]{Kollar2023Families}.
A normal variety is $\mathbb Q$-factorial if every prime Weil divisor has a
positive Cartier multiple.  A divisor is nef over a morphism when it has
nonnegative degree on every complete contracted curve.
